\documentclass[10pt,a4paper]{amsart}
\usepackage[margin=19mm]{geometry}
\usepackage{amsmath,amssymb,mathtools}
\usepackage[scr=rsfs,cal=euler]{mathalfa}
\usepackage{xfrac}
\usepackage[unicode,hidelinks,bookmarks=false]{hyperref}
\newcommand{\ie}{i.e.\ }
\newcommand{\cf}{cf.\ }
\newcommand{\resp}{respectively\ }
\newcommand{\Subsection}{\subsection}
\providecommand{\nobreakdash}{\nobreak\hbox{-}}
\setlength{\emergencystretch}{2em}
\multlinegap0pt
\usepackage{bm}
\usepackage{algorithm}\usepackage{algpseudocode}
\usepackage{xcolor}
\usepackage{csquotes}
\usepackage{float}
\usepackage{amssymb}
\newtheorem{prop}{Proposition}
\title{Next Generation Equation-Free Multiscale Modelling of Crowd Dynamics via Machine Learning}
\author{Hector Vargas Alvarez, Dimitrios G. Patsatzis, Lucia Russo, Ioannis Kevrekidis, Constantinos Siettos}
\date{Source: arXiv:2508.03926v3 (CC BY 4.0)}
\begin{document}
\maketitle

\noindent\emph{Source and licence.} Next Generation Equation-Free Multiscale Modelling of Crowd Dynamics via Machine Learning. Version arXiv:2508.03926v3 (CC BY 4.0), distributed by arXiv under the Creative Commons Attribution 4.0 International licence, http://creativecommons.org/licenses/by/4.0/, as recorded in the arXiv licence metadata for this exact version and shown on the abstract page https://arxiv.org/abs/2508.03926v3. Full licence notice: \texttt{licenses/arxiv-2508.03926-CC-BY-4.0.txt}.\par\smallskip

\noindent\textit{Editorial scope.} Counted unit: the article's prop prop:2 (source0.tex lines 308--320). The article's complete original proof of it is delivered with it (source0.tex lines 321--380, one \texttt{proof} environment of 60 lines). No mathematics is written, repaired or re-derived: the statement and the proof are the article's own lines, in the article's own notation, and no retained line is substituted or re-flowed.  Retained uncounted as the source-local context the proof needs: lines 275--278; lines 283--285.  Nothing was omitted from any retained span, so the package carries no omission record.  No retained line contains a citation, so the wrapper carries no bibliography and no bibliography entry.  No retained line contains a figure environment, an image-inclusion command or drawing code, so no image is delivered and the package declares no asset dependency.  Editorial formatting only, in addition: an AMS \texttt{amsart} preamble that restates the article's own declarations and macros used by the retained lines, namely the environments \texttt{prop} on separate counters, together with the standard packages those lines need.  The retained environments print in this document as Proposition 1 in order of appearance, whereas the article prints the same items with the numbering of its own full document.  The complete native source of the article as distributed in the arXiv e-print for this exact version is bundled unchanged as \texttt{sources/182766/source0.tex}.
\begin{equation}
   \mathbf{1}_{n_c}^\top \bar{X}= \mathbf{1}_{n_c}^\top XH=\mathbf{1}_{n_t}^\top H=\mathbf{1}_{n_t}^\top-\frac{1}{n_t}\mathbf{1}_{n_t}^\top\mathbf{1}_{n_t}\mathbf{1}_{n_t}^\top=\mathbf{1}_{n_t}^\top-\mathbf{1}_{n_t}^\top=\boldsymbol{0}^\top_{n_t}.
   \label{eq:orthog1}
\end{equation}

\begin{equation} \label{eq:ortho1a}
    \mathbf{1}_{n_c}^\top U_d = \boldsymbol{0}^\top_d,
\end{equation}

\begin{prop} \label{prop:2}
Let us assume the snapshot matrices $X^{(1)},\, X^{(2)} \in\mathbb{R}^{n_c \times n_t}$ with columns being normalized density fields of two interacting population groups. Furthermore, let $W\in \mathbb{R}^{n_c \times m}$ and $T \in \mathbb{R}^{n_c \times m}$ be the left and right singular vectors of the cross-covariance matrix $C^{(12)}=\frac{1}{n_t}\bar X^{(1)} \left(\bar X^{(2)}\right)^\top \in \mathbb{R}^{n_c\times n_c}$, and consider the augmented orthonormal projection bases: 
\begin{equation} \label{eq:Aug_Basis}
    \tilde U^{(1)} = \begin{bmatrix} \frac{1}{\sqrt{n_c}}\mathbf{1}_{n_c} & U_{d_1}^{(1)} & \hat W \end{bmatrix} \in \mathbb{R}^{n_c \times (d_1+m+1)}, \quad
\tilde U^{(2)} = \begin{bmatrix} \frac{1}{\sqrt{n_c}}\mathbf{1}_{n_c} & U_{d_2}^{(2)} & \hat T \end{bmatrix} \in \mathbb{R}^{n_c \times (d_2+m+1)},
\end{equation} 
where $U_{d_1}^{(1)} \in \mathbb{R}^{n_c\times d_1}$ and $U_{d_2}^{(2)} \in \mathbb{R}^{n_c\times d_2}$ are the first $d_1$ and $d_2$ left singular vectors from the SVD of $\bar X^{(1)}$ and $\bar X^{(2)}$, respectively, and the terms $\hat W = W_\perp (W_\perp^\top W_\perp)^{-1/2}$ and $\hat T = T_\perp (T_\perp^\top T_\perp)^{-1/2}$ are computed by the orthogonal projections $W_\perp = \left(I_{n_c} - \mathbf{1}_{n_c}\mathbf{1}^{\top}_{n_c}/n_c - U_{d_1}^{(1)} (U_{d_1}^{(1)})^\top\right)W \in \mathbb{R}^{n_c\times m}$ and $T_\perp = \left(I_{n_c} - \mathbf{1}_{n_c}\mathbf{1}^{\top}_{n_c}/n_c- U_{d_2}^{(2)} (U_{d_2}^{(2)})^\top\right)T \in \mathbb{R}^{n_c\times m}$. If the original densities are mass-preserving along time steps, i.e., if $\mathbf{1}^{\top}_{n_c}X^{(1)}=\mathbf{1}^{\top}_{n_{t}}$ and $\mathbf{1}^{\top}_{n_c}X^{(2)}=\mathbf{1}^{\top}_{n_{t}}$, then the reconstructed densities computed by projecting $X^{(1)}$ and $X^{(2)}$ onto the augmented bases as: 
\begin{equation}
    \tilde X^{(1)}=\tilde {U}^{(1)} \left(\tilde {U}^{(1)}\right)^{\top} \bar X^{(1)}+X^{(1)} (I_{n_t}-H), \quad \tilde X^{(2)}=\tilde {U}^{(2)} \left(\tilde {U}^{(2)}\right)^{\top} \bar X^{(2)}+X^{(2)} (I_{n_t}-H),
    \label{Eq:Reconstruction_agumented}
\end{equation}
are also mass-preserving, i.e., $\mathbf{1}^{\top}_{n_c}\tilde X^{(1)}=\mathbf{1}^{\top}_{n_{t}}$ and $\mathbf{1}^{\top}_{n_c} \tilde X^{(2)}=\mathbf{1}^{\top}_{n_{t}}$.
\end{prop}

\begin{proof}
Assume the original densities are mass-preserving for each group. From the properties of the single-population POD (see Eq.~\eqref{eq:orthog1} and Eq.~\eqref{eq:ortho1a}), it follows that:
\begin{equation} \label{eq:singlePop_ortho}
   \mathbf{1}_{n_c}^\top \bar{X}^{(l)}=\boldsymbol{0}^\top_{n_t}, \quad \mathbf{1}_{n_c}^\top U^{(l)}_{d_l} = \mathbf{0}_{d_l}^\top, 
\end{equation}
where $\bar{X}^{(l)}=X^{(l)}H$ is the centered data matrix and $U_{d_l}^{(l)}\in \mathbb{R}^{n_c\times d_l}$ are the first $d_l$ left singular vectors from the SVD of $\bar X^{(l)}$ for the $l=1,2$ groups. Let's consider now the SVD of the cross-covariance matrix: 
\begin{equation}
    C^{(12)} = \frac{1}{n_t}\bar X^{(1)} \left(\bar X^{(2)}\right)^\top= W \Sigma_{C^{(12)}} T^{\top}, %\quad  W \in \mathbb{R}^{n_{c1} \times k}, \quad \Sigma_{C^{(12)}} \in \mathbb{R}^{m \times m}, \quad T \in \mathbb{R}^{n_{c2} \times k},
    \label{Eq:Cross-cov_SVD}
\end{equation}
with $W \in \mathbb{R}^{n_c \times m}$, $\Sigma_{C^{(12)}} \in \mathbb{R}^{m \times m}$ and $T \in \mathbb{R}^{n_c \times m}$, where $m$ denotes the number of retained cross-covariance modes. Each column of $C^{(12)}$ is a linear combination of the columns of $\bar X^{(1)}$, and each row a linear combination of the rows of $\bar X^{(2)}$. Consequently:
\begin{equation}
\mathcal{R}(W)\equiv \mathcal{R}\big(C^{(12)}\big) \subseteq \mathcal{R}(\bar{X}^{(1)}), %\equiv\mathcal{R}\big(U^{(1)}_{d_1}\big), 
\quad 
\mathcal{R}(T)\equiv\mathcal{R}\big((C^{(12)})^\top\big) \subseteq \mathcal{R}\big(\bar{X}^{(2)}\big). %\equiv \mathcal{R}\big(U^{(2)}_{d_2}\big).
\end{equation}
%implying that the left singular vectors $W$ of $C^{(12)}$, are related to $\bar{X}^{(1)}$ and the right singular vectors $T$ are related to $\bar{X}^{(2)}$.

To obtain interaction modes orthogonal to both the unity vectors $\mathbf{1}_{n_c}$ and the individual POD bases $U^{(l)}_{d_l}$ for $l=1,2$, we project the $W$ and $T$ onto the complementary subspace: 
\begin{equation}
    W_\perp = \left(I_{n_c} - \frac{1}{n_c}\mathbf{1}_{n_c}\mathbf{1}^{\top}_{n_c} - U_{d_1}^{(1)} (U_{d_1}^{(1)})^\top\right)W \in \mathbb{R}^{n_c\times m}, \quad T_\perp = \left(I_{n_c} - \frac{1}{n_c} \mathbf{1}_{n_c}\mathbf{1}^{\top}_{n_c}- U_{d_2}^{(2)} (U_{d_2}^{(2)})^\top\right)T \in \mathbb{R}^{n_c\times m},
\end{equation}
and define their orthonormalized versions as:
\begin{equation}
    \hat W = W_\perp (W_\perp^\top W
    _\perp)^{-1/2}, \quad  \hat T = T_\perp (T_\perp^\top T_\perp)^{-1/2}.
\end{equation}
By construction of $\hat{W}$ and $\hat{T}$, we obtain:
\begin{equation} \label{eq:cross_ortho}
    \mathbf{1}_{n_c}^\top \hat{W}= \mathbf{0}^\top_m, \quad (U_{d_1}^{(1)})^\top \hat{W}=0\in \mathbb{R}^{d_1 \times n_c}, \quad
    \mathbf{1}_{n_c}^\top \hat{T}= \mathbf{0}^\top_k, \quad (U_{d_2}^{(2)})^\top \hat{T}=0\in \mathbb{R}^{d_2 \times n_c}.
\end{equation}
The augmented orthonormal bases (Eq.~\eqref{eq:Aug_Basis}) are therefore:
\begin{equation} \label{eq:Aug_Basis_proof}
    \tilde U^{(1)} = \begin{bmatrix} \frac{1}{\sqrt{n_c}}\mathbf{1}_{n_c} & U_{d_1}^{(1)} & \hat W \end{bmatrix} \in \mathbb{R}^{n_c \times (d_1+m+1)}, \quad
\tilde U^{(2)} = \begin{bmatrix} \frac{1}{\sqrt{n_c}}\mathbf{1}_{n_c} & U_{d_2}^{(2)} & \hat T \end{bmatrix} \in \mathbb{R}^{n_c \times (d_2+m+1)},
\end{equation} 
according to which, the reduced coordinates of the snapshots can be obtained as:
\begin{equation}
Y^{(1)} = (\tilde{U}^{(1)})^\top \bar X^{(1)} \in \mathbb{R}^{(d_1+m+1) \times n_t}, \qquad
Y^{(2)} = (\tilde{U}^{(2)})^\top \bar X^{(2)} \in \mathbb{R}^{(d_2+m+1) \times n_t}.
\end{equation}

Multiplying the augmented basis in Eq.~\eqref{eq:Aug_Basis_proof} by $\mathbf{1}_{n_c}^\top$ from the left, and using Eq.~\eqref{eq:cross_ortho} together with the orthogonality $\mathbf{1}_{n_c}^\top U^{(l)}_{d_l} = \mathbf{0}_{d_l}^\top$ in Eq.~\eqref{eq:singlePop_ortho} gives:
\begin{equation}
    \mathbf{1}_{n_c}^\top \tilde U^{(1)} = \sqrt{n_c} \mathbf{e}_1^\top, \quad
    \mathbf{1}_{n_c}^\top \tilde U^{(2)} = \sqrt{n_c} \mathbf{e}_1^\top,
\end{equation}
where $\mathbf{e}_1=[1,0,\ldots,0]^\top$. Hence, for the projected centered data, we get:
\begin{equation} \label{eq:cross_ortho1}
    \mathbf{1}_{n_c}^\top \tilde U^{(l)} (\tilde{U}^{(l)})^\top \bar X^{(l)} = \sqrt{n_c} \mathbf{e}_1^\top (\tilde{U}^{(l)})^\top \bar X^{(l)} = \mathbf{1}_{n_c}^\top \bar X^{(l)} = \mathbf{0}_{n_t}^\top,\quad l=1,2.
\end{equation}
The last equality uses the mass-preservation property for the original densities $ \mathbf{1}_{n_c}^\top \bar{X}^{(l)}=\boldsymbol{0}^\top_{n_t}$ in Eq.~\eqref{eq:singlePop_ortho}.

Finally, using the reconstruction formula in Eq.~\eqref{Eq:Reconstruction_agumented}, and the identities derived in Eqs.~\eqref{eq:singlePop_ortho} and \eqref{eq:cross_ortho1}, we obtain:
\begin{align}
\mathbf{1}^{\top}_{n_c}\tilde X^{(l)}=\mathbf{1}_{n_c}^\top \tilde U^{(l)} (\tilde{U}^{(l)})^\top \bar X^{(l)}+\mathbf{1}^{\top}_{n_c}  X^{(l)} (I_{n_t}-H)=\mathbf{1}^{\top}_{n_c}  X^{(l)}(I_{n_t}-H)=\mathbf{1}_{n_t},\quad l=1,2. 
\end{align}
Thus, the reconstructions $\tilde X^{(1)}$ and $\tilde X^{(2)}$, obtained via the projection of $X^{(1)}$ and $X^{(2)}$ onto the augmented basis in Eq.~\eqref{eq:Aug_Basis_proof}, retain the mass-preservation property of the original data.
\end{proof}

\end{document}
